\bookchapter{Capitolo 36: Telefono e casa intelligente — rendi la casa smart}

\chapterintro{Il telefono può fare una cosa sempre più importante: diventare il «telecomando» dell'intera casa. Da una lampadina a una presa, da una serratura a un sensore, la casa intelligente sta passando da «giocattolo da smanettoni» a «dotazione di casa». Questo capitolo ti aiuta a fare chiarezza su protocolli, ecosistemi, hub e kit di partenza, e si chiude con una checklist di sicurezza per non farti mordere dallo «smart».}

\booksection{Protocolli della casa intelligente: la scelta conta}

Prima di comprare un dispositivo smart controlla quale protocollo supporta: determina se può essere «comandato» dal tuo telefono e dal tuo ecosistema. I principali sono quattro, con caratteristiche diverse:

\begin{itemize}
  \item Protocollo Matter: pubblicato nel 2022 dalla CSA (Connectivity Standards Alliance) con dietro Apple, Google, Amazon, Samsung e altri giganti. Unico obiettivo: far parlare tutti i dispositivi smart la stessa lingua. Nel 2023-2026 sta diventando rapidamente lo standard unificato; i dispositivi compatibili si collegano a Apple HomeKit, Google Home, Amazon Alexa e Mi Home senza restare legati a un solo ecosistema.
  \item Protocollo Thread: il «tubature» di rete alla base di Matter. Basso consumo, rete mesh (se un nodo cade, il segnale fa un'altra strada), risposta rapida e senza scatti. I dispositivi con router di bordo Thread (HomePod mini, Apple TV 4K, gateway Xiaomi) danno un'esperienza nettamente più stabile.
  \item Wi-Fi: banda ampia, connessione diretta al router, ma alto consumo e con molti dispositivi si intasa il canale; adatto a telecamere e robot aspirapolvere che trasferiscono molti dati.
  \item Bluetooth (BLE): configurazione semplice e risparmio energetico, ma breve distanza e niente mesh; adatto a serrature, termometri e igrometri.
  \item Zigbee: protocollo mesh a basso consumo storico, richiede un gateway dedicato; dispositivi economici e stabili ma scarsa interoperabilità tra marchi; sta progressivamente lasciando il posto a Matter+Thread.
\end{itemize}

\begin{tcbinfo}{Principio in una frase}
Nel 2026, per i dispositivi smart guarda prima il marchio «Matter over Thread»: è la scelta più resistente al tempo. I vecchi dispositivi solo Wi-Fi o Bluetooth non vanno buttati subito, usali e ai prossimi acquisti passa a Matter.

\end{tcbinfo}

\booksection{Confronto tra i tre grandi ecosistemi}

Il protocollo è la «lingua», l'ecosistema è il «circolo». Anche se tutto supporta Matter, nella vita quotidiana devi scegliere un ecosistema principale per automazioni e controllo vocale. Confronto tra i tre principali:

\begin{center}
\begin{tabularx}{\linewidth}{|X|X|X|X|}
\hline
\textbf{Ecosistema} & \textbf{Punti di forza} & \textbf{Punti deboli} & \textbf{Per chi} \tabularnewline
\hline
Apple HomeKit & Privacy massima (elaborazione locale), integrazione fluida con iPhone/Siri/Scorciatoie & Pochi dispositivi e costosi, serve HomePod o Apple TV come hub & Famiglie con dispositivi Apple, attenti a privacy ed esperienza \tabularnewline
Mi Home (Xiaomi) & Ampia gamma di dispositivi, ottimo rapporto qualità-prezzo, Xiao AI comoda & Ecosistema relativamente chiuso, molti dati su cloud, molti annunci & Chi ha budget limitato e vuole attrezzare tutta casa in una volta \tabularnewline
Huawei HiLink (HarmonyOS) & Continuità perfetta con telefono/tablet/orologio HarmonyOS, gamma di dispositivi in rapida crescita & Pochi dispositivi supportati all'estero, ecosistema ancora in recupero & Utenti Huawei, attenti alla collaborazione multi-dispositivo «Super Device» \tabularnewline
\hline
\end{tabularx}
\end{center}

\begin{tcbinfo}{Consigli per scegliere l'ecosistema}
Scegli in base al telefono: utenti iPhone preferiscano HomeKit, utenti Huawei HiLink, altri Android Mi Home è il più pratico. Grazie a Matter, dispositivi di ecosistemi diversi possono essere controllati dallo stesso telefono: non devi schierarti subito.

\end{tcbinfo}

\booksection{Il telefono come hub di casa}

Nella maggior parte dei casi iniziali il telefono è il tuo «hub di casa»: con l'app controlli luci, condizionatori, tende e sicurezza, e crei scenari automatici per far «muovere» la casa da sola:

\begin{itemize}
  \item Controllo dispositivi: con HomeKit/Mi Home/Smart Life apri e spegni le luci da remoto o in locale, regola la temperatura colore, imposta l'aria condizionata, apri e chiudi le tende, guarda le telecamere.
  \item Scenari automatici: combina «condizione azione»; i tre più usati: Modalità rientro: arrivo a casa + buio accendi luce ingresso + aria condizionata + chiudi tende + chiudi serratura.
  \item Modalità uscita: l'ultimo che esce (posizione fuori casa o serratura chiusa) spegni luci + aria + prese + attiva telecamere + allarme.
  \item Modalità notte: ogni sera alle 23:00 spegni luci principali + chiudi tende + aria a 26 gradi + attiva lampada da notte.
\end{itemize}

Tocco NFC: attacca un tag NFC alla porta e con un tocco del telefono attivi lo scenario (un tocco luce ingresso, un altro modalità notte). Su iPhone con le «Scorciatoie», su Android con «Wallet» o l'app del produttore.

\begin{tcbpractice}{Prova: crea una «Modalità rientro»}
Nella sezione «Automazioni» dell'app casa intelligente creane una nuova: condizione «Arrivo a un luogo» (casa), azioni in sequenza «accendi luce ingresso» «aria a 26 gradi» «chiudi tende». Salva, esci a fare una passeggiata e al rientro guarda se le luci si accendono da sole: la soddisfazione della prima volta è meglio di comprare un nuovo dispositivo.

\end{tcbpractice}

\booksection{Kit di partenza consigliati}

Non serve attrezzare tutto subito: con questi quattro elementi copri l'80\% delle esigenze quotidiane:

\begin{itemize}
  \item Altoparlante smart: Xiao AI (Mi Home, ottimo rapporto qualità-prezzo) / HomePod mini (Apple, funge da hub Thread) / Alexa (ampio ecosistema estero). Attenzione: l'altoparlante determina l'ecosistema vocale, scegli prima l'ecosistema.
  \item Lampadina smart: Yeelight (Mi Home/Matter, temperatura e luminosità regolabili, economica) / Philips Hue (colori belli e stabili, ma serve il gateway Hue, più costosa). Attenzione: controlla la base (E27/B22) e se serve la regolazione colore.
  \item Presa smart: Xiaomi Smart Plug (economica, contatore consumi) / TP-Link Kasa (ottimo supporto Matter all'estero). Attenzione: con elettrodomestici potenti (scaldabagno, stufa) controlla la potenza nominale per non sovraccaricare.
  \item Serratura intelligente: Loock/Xiaomi (rapporto qualità-prezzo, integrazione Mi Home) / Kaadas (produttore tradizionale, struttura meccanica solida). Attenzione: verifica il supporto Matter o del tuo ecosistema, l'autonomia della batteria e la posizione della serratura di emergenza.
\end{itemize}

\begin{tcbinfo}{Ordine di acquisto consigliato}
Prima l'altoparlante (fissa l'ecosistema) poi lampadina e presa (facili, feedback immediato) infine serratura e telecamere (sicurezza, scegli con cura). Aggiungi un tipo di dispositivo alla volta, famigliarizza e poi espandi, per non ritrovarti con cose che non sai usare.

\end{tcbinfo}

\booksection{Attenzione alla sicurezza}

Più dispositivi smart, più grande la superficie di attacco. Queste quattro regole sono il minimo sindacale, controllale una per una:

\begin{enumerate}
  \item Cambia le password predefinite: le password di fabbrica di telecamere, router e serrature sono pubbliche online; la prima cosa da fare è impostare una password «lunga e non ripetuta».
  \item Aggiorna il firmware: gli aggiornamenti dei produttori spesso correggono vulnerabilità; attiva «Aggiornamento automatico» nell'app o almeno controlla una volta al mese.
  \item Isolamento di rete: metti i dispositivi IoT su un Wi-Fi separato (o rete ospiti), distinto dalla rete principale di telefono e computer. Se un dispositivo viene compromesso, è più difficile che minacci lateralmente i tuoi dati importanti.
  \item Protezione della privacy: sulle telecamere attiva la mascheratura delle «aree private» (camera da letto, bagno non ripresi), disattiva la registrazione cloud 24 ore; non lasciare attive a lungo le telecamere di casa ad app di terze parti.
\end{enumerate}

\textbf{Errori da non commettere}\par

\begin{itemize}
  \item Non usare l'«apertura da remoto» della serratura per sconosciuti: preferisci password temporanee che scadono dopo l'uso.
  \item Non puntare la telecamera verso il letto o fuori dalla finestra: stai esponendo la tua privacy.
  \item Non comprare per risparmiare dispositivi senza marca: quasi mai aggiornano il firmware e sono una porta aperta in casa.
\end{itemize}

\begin{tcbsummary}{Riepilogo del capitolo}
\begin{itemize}
  \item Matter+Thread è la direzione unificata 2023-2026: per i nuovi dispositivi guarda prima questo marchio
  \item I tre ecosistemi hanno punti di forza diversi: Apple privacy ed esperienza, Mi Home rapporto qualità-prezzo, Huawei collaborazione multi-dispositivo
  \item Il telefono può fare da hub; con gli scenari «rientro/uscita/notte» la casa si muove da sola
  \item Kit iniziale: altoparlante, lampadina, presa, serratura, in quest'ordine
  \item Password predefinite da cambiare, firmware aggiornati, rete isolata, privacy protetta: quattro regole di sicurezza da non trascurare
\end{itemize}

\end{tcbsummary}
